\documentclass{skripsi}
% generate minimum example for testing dtx code snippets
\title{Prediksi Kerusakan Jalan Raya dengan Machine Learning}
\author{Ipsum}
\authorid{1234567890}
\date{\today}
\faculty{Fakultas Teknik}
\program{Program Studi Teknik Sipil}
\university{Universitas Muhammadiyah Yogyakarta}
\yearofsubmission{\the\year}
\PassOptionsToPackage{style=apa}{biblatex}
\usepackage{lipsum}
\begin{document} 
% \maketitle
\frontmatter
\onehalfspacing
\tableofcontents
\listoffigures
\listoftables
\mainmatter
\pagestyle{fancyplain}
\onehalfspacing
\chapter{Lorem}
\section{Consectetur}
Lorem ipsum dolor sit asmet, consectetur adipiscing elit. Sed non risus. Suspendisse lectus tortor, dignissim sit amet, adipiscing nec, ultricies sed, dolor. Cras elementum ultrices diam. Maecenas ligula massa, varius a, semper congue, euismod non, mi. Proin porttitor, orci nec nonummy molestie, enim est eleifend mi, non fermentum diam nisl sit amet erat. Duis semper. Duis arcu massa, scelerisque vitae, consequat in, pretium a, enim. Pellentesque congue. Ut in risus volutpat libero pharetra tempor. Cras vestibulum bibendum augue. Praesent egestas leo in pede. Praesent blandit odio eu enim. Pellentesque sed dui ut augue blandit sodales. Vestibulum ante ipsum primis in faucibus orci luctus et ultrices posuere cubilia Curae; Aliquam nibh. Mauris ac mauris sed pede pellentesque fermentum. Maecenas adipiscing ante non diam sodales hendrerit.
% create lorem ipsum using package with subsections like a thesis up to chapter 5 with some dummy figures example-image-b and tables
\section{Dolor}
\lipsum[1-3]
\subsection{Sit Amet}
\lipsum[4-6]
\subsubsection{Consectetur}
\lipsum[7-9]

\chapter{Ipsum}
% give an example of a figure
\section{Amet}
\begin{figure}[h]
    \centering
    \includegraphics[width=0.5\textwidth]{example-image-b}
    \caption{Contoh Gambar}
    \label{fig:contoh-gambar}
\end{figure}
\lipsum[10-12]
% give an example of a table
\begin{table}[bh]
    \centering
    \begin{tabular}{|c|c|c|}
        \hline
        Kolom 1 & Kolom 2 & Kolom 3 \\
        \hline
        Data 1 & Data 2 & Data 3 \\
        Data 4 & Data 5 & Data 6 \\ 
        \hline
    \end{tabular}
    \caption{Contoh Tabel}
    \label{tab:contoh-tabel}
\end{table}

\end{document}